\begin{proposition}\label{sga3:XXI:3.1.2}
Let \(\alpha\) and \(\alpha_i\), \(i=1,\ldots,n\), be roots. If
\[
 \alpha=\sum_{i=1}^n m_i\alpha_i,
 \qquad m_i\in\mathbb N,
\]
then there is a sequence \(\beta_1,\ldots,\beta_m\) of roots chosen
from among the \(\alpha_i\), not necessarily pairwise distinct, such
that, on setting
\[
 \gamma_p=\sum_{i=1}^p\beta_i,\qquad p=1,\ldots,m,
\]
one has \(\gamma_p\in R\) and \(\gamma_m=\alpha\).
\end{proposition}

\begin{proof}
We argue by induction on the integer
\(m'=\sum_i m_i\). If \(\alpha\) equals one of the \(\alpha_i\), say
\(\alpha_{i_0}\), which necessarily happens when \(m'=1\), take
\[
 m=1,\qquad\beta_1=\alpha_{i_0}.
\]
Otherwise, Lemma~\ref{sga3:XXI:3.1.1} gives an index \(i\) for which
\(m_i\neq0\) and \(\alpha-\alpha_i\) is a root. Then
\(m_i-1\in\mathbb N\) and
\[
 \alpha-\alpha_i
 =(m_i-1)\alpha_i+\sum_{j\neq i}m_j\alpha_j.
\]
Apply the induction hypothesis to \(\alpha-\alpha_i\), and append
\(\alpha_i\) to the resulting sequence.
\end{proof}

\begin{corollary}\label{sga3:XXI:3.1.3}
Let \(R'\subset R\). The following conditions are equivalent:
\begin{enumerate}
\item[\SGAInlineRefTarget{sga3:XXI:3.1.3:item:i}(i)] if \(\alpha,\beta\in R'\) and \(\alpha+\beta\in R\), then
  \(\alpha+\beta\in R'\);
\item[\SGAInlineRefTarget{sga3:XXI:3.1.3:item:ii}(ii)] \((\mathbb N\cdot R')\cap R\subset R'\).
\end{enumerate}
\end{corollary}

\begin{proof}
Clearly \SGARefLink{sga3:XXI:3.1.3:item:ii}{(ii)} implies \SGARefLink{sga3:XXI:3.1.3:item:i}{(i)}. The converse follows immediately from
Proposition~\ref{sga3:XXI:3.1.2}.
\end{proof}

\begin{definition}\label{sga3:XXI:3.1.4}
A set of roots satisfying the equivalent conditions of
Corollary~\ref{sga3:XXI:3.1.3} is called \emph{closed}.
\end{definition}

\begin{proposition}\label{sga3:XXI:3.1.5}
Let \(\Delta\subset R\) be a set of roots. The following conditions
are equivalent:
\begin{enumerate}
\item[\SGAInlineRefTarget{sga3:XXI:3.1.5:item:i}(i)] the elements of \(\Delta\) are indivisible and linearly
  independent, and
  \[
   R\subset(\Qp\cdot\Delta)\cup(-\Qp\cdot\Delta);
  \]
\item[\SGAInlineRefTarget{sga3:XXI:3.1.5:item:ii}(ii)] the elements of \(\Delta\) are linearly independent, and
  \[
   R\subset(\mathbb N\cdot\Delta)\cup(-\mathbb N\cdot\Delta);
  \]
\item[\SGAInlineRefTarget{sga3:XXI:3.1.5:item:iii}(iii)] every root can be written uniquely as a linear
  combination of the elements of \(\Delta\), with integer
  coefficients all of the same sign.
\end{enumerate}
\end{proposition}

\begin{proof}
\footnote{Editors' note: in the argument that follows, the order was
changed and the proof of the implication \SGARefLink{sga3:XXI:3.1.5:item:iii}{(iii)}\(\Rightarrow\)\SGARefLink{sga3:XXI:3.1.5:item:i}{(i)} was
spelled out.}
It is clear that \SGARefLink{sga3:XXI:3.1.5:item:ii}{(ii)} implies \SGARefLink{sga3:XXI:3.1.5:item:iii}{(iii)}. Write
\(\alpha_1,\ldots,\alpha_n\) for the distinct elements of \(\Delta\).

Assume \SGARefLink{sga3:XXI:3.1.5:item:i}{(i)}, and let \(\alpha\in R\). There is a unique expression
\[
 \varepsilon\alpha=\sum_i q_i\alpha_i,
 \qquad q_i\in\Qp,\qquad\varepsilon=\pm1.
\]
We show that the \(q_i\) are integers. This is immediate when all but
one of them vanish, by Corollary~\ref{sga3:XXI:2.1.5}. Otherwise,
\(\varepsilon\alpha\) is distinct from the \(\alpha_i\), and
Lemma~\ref{sga3:XXI:3.1.1} gives an \(i_0\) such that
\[
 \alpha'=\varepsilon\alpha-\alpha_{i_0}\in R,
 \qquad q_{i_0}\neq0.
\]
Thus
\[
 \alpha'
 =(q_{i_0}-1)\alpha_{i_0}
  +\sum_{i\neq i_0}q_i\alpha_i.
\]
At least one \(q_i\), with \(i\neq i_0\), is nonzero, so \SGARefLink{sga3:XXI:3.1.5:item:i}{(i)} forces
\(q_{i_0}-1\geq0\). Repeating this operation with \(\alpha'\), after
finitely many steps we find that all \(q_i\) are integers. This proves
\SGARefLink{sga3:XXI:3.1.5:item:ii}{(ii)}.

Assume now \SGARefLink{sga3:XXI:3.1.5:item:iii}{(iii)}. We first prove that
\(\alpha_1,\ldots,\alpha_n\) are linearly independent over
\(\mathbb Q\). Otherwise there would be an equality
\[
 x=\sum_{i\in I}a_i\alpha_i
  =\sum_{j\in J}b_j\alpha_j,
\]
where \(I\) and \(J\) are disjoint subsets of
\(\{1,\ldots,n\}\), at least one of them, say \(I\), is nonempty, and
\(a_i,b_j\in\mathbb N^\ast\). By the corollary in
Complements~\ref{sga3:XXI:2.3.7}, one has
\(\langle\alpha_i,\alpha_j\rangle\leq0\) for \(i\in I\) and \(j\in J\).
Therefore \(\langle x,x\rangle\leq0\), and hence \(x=0\).

Choose \(i_0\in I\). The two expressions
\[
 \alpha_{i_0}
 =\alpha_{i_0}+\sum_{i\in I}a_i\alpha_i
 =\alpha_{i_0}+\sum_{j\in J}b_j\alpha_j
\]
then imply, by \SGARefLink{sga3:XXI:3.1.5:item:iii}{(iii)}, that \(a_i=0=b_j\) for all \(i,j\), a
contradiction. Thus the elements of \(\Delta\) are linearly
independent over \(\mathbb Q\).

It remains to show that they are indivisible. Suppose that
\(\beta\in\Delta\) is divisible. Then \(\beta/2\in R\), so
\[
 \frac{\beta}{2}
 =\varepsilon\sum_{\alpha\in\Delta}m_\alpha\alpha,
 \qquad m_\alpha\in\mathbb N,\qquad\varepsilon=\pm1.
\]
Consequently,
\[
 \beta=2\varepsilon\sum_{\alpha\in\Delta}m_\alpha\alpha.
\]
Uniqueness gives \(m_\alpha=0\) for \(\alpha\neq\beta\) and
\(2\varepsilon m_\beta=1\), a contradiction.
\end{proof}

\begin{definition}\label{sga3:XXI:3.1.6}
A set \(\Delta\) of roots satisfying the conditions of
Proposition~\ref{sga3:XXI:3.1.5} is called a
\emph{system of simple roots}, or a \emph{base} of \(\cR\).
\end{definition}

If \(w\in W\) and \(\Delta\) is a system of simple roots, then
\(w(\Delta)\) is a system of simple roots.

\begin{remark}\label{sga3:XXI:3.1.7}
This definition involves only \(R\), not \(\cR\); indeed, it involves
only \(\operatorname{ind}(R)\).
\end{remark}

\begin{remark}\label{sga3:XXI:3.1.8}
If \(\Delta\) is a system of simple roots, then \(\Delta\) is a basis
of the free abelian group \(\Gamma_0(R)\). Hence
\[
 \operatorname{Card}(\Delta)=\rgss(\cR).
\]
\end{remark}

\begin{remark}\label{sga3:XXI:3.1.9}
The conditions of Proposition~\ref{sga3:XXI:3.1.5} are therefore
equivalent to the following:
\begin{enumerate}
\item[(i\('\))] the elements of \(\Delta\) are indivisible, their
  number is at most \(\rgss(\cR)\), and
  \[
   R\subset(\Qp\cdot\Delta)\cup(-\Qp\cdot\Delta);
  \]
\item[(ii\('\))] \(\operatorname{Card}(\Delta)\leq\rgss(\cR)\), and
  \[
   R\subset(\mathbb N\cdot\Delta)\cup(-\mathbb N\cdot\Delta).
  \]
\end{enumerate}
\end{remark}
